\documentclass[a4paper,12pt]{article}
\usepackage[utf8]{inputenc}
\usepackage[T1]{fontenc}
\usepackage[ngerman]{babel}
\usepackage{amsmath}
\usepackage{amssymb}
\usepackage{enumitem}
\usepackage{geometry}
\geometry{a4paper, margin=2.5cm}

\title{Einsendeaufgaben -- Lektion 3\\[0.5em]
\large Modul 61111: Mathematische Grundlagen}
\author{}
\date{}

\begin{document}
\maketitle

\section*{Aufgabe 3.2}

\begin{enumerate}[label=\alph*)]

    \item \textbf{Bestimmung Basis von $U \cap V$:}

    \[
        U \cap V = \left\{ v \in \mathbb{R}^3 \;\middle|\;
        \exists\, a_1, a_2, b_1, b_2 :\;
        a_1 \begin{pmatrix}1\\2\\3\end{pmatrix}
        + a_2 \begin{pmatrix}-1\\2\\1\end{pmatrix}
        = b_1 \begin{pmatrix}-1\\0\\1\end{pmatrix}
        + b_2 \begin{pmatrix}0\\-1\\1\end{pmatrix}
        \right\}
    \]

    Also:
    \[
        b_1 \begin{pmatrix}1\\0\\-1\end{pmatrix}
        + b_2 \begin{pmatrix}0\\1\\-1\end{pmatrix}
        + a_1 \begin{pmatrix}1\\2\\3\end{pmatrix}
        + a_2 \begin{pmatrix}-1\\2\\1\end{pmatrix}
        = \mathbf{0}
    \]

    \[
        \Longleftrightarrow
        \begin{pmatrix}1 & 0 & 1 & -1\\ 0 & 1 & 2 & 2\\ -1 & -1 & 3 & 1\end{pmatrix}
        \begin{pmatrix}b_1\\b_2\\a_1\\a_2\end{pmatrix} = \mathbf{0}
    \]

    Lösen der homogenen Gleichung:
    \[
        \begin{pmatrix}1 & 0 & 1 & -1\\ 0 & 1 & 2 & 2\\ -1 & -1 & 3 & 1\end{pmatrix}
        \;\Big|\; Z_{3,1}(1)
        \longrightarrow
        \begin{pmatrix}1 & 0 & 1 & -1\\ 0 & 1 & 2 & 2\\ 0 & -1 & 4 & 0\end{pmatrix}
        \;\Big|\; Z_{3,2}(1)
    \]
    \[
        \longrightarrow
        \begin{pmatrix}1 & 0 & 1 & -1\\ 0 & 1 & 2 & 2\\ 0 & 0 & 6 & 2\end{pmatrix}
        \;\Big|\;
        \begin{array}{l}
            Z_{1,3}\!\left(-\tfrac{1}{6}\right)\\[2pt]
            Z_{2,3}\!\left(-\tfrac{1}{3}\right)\\[2pt]
            Z_3\!\left(\tfrac{1}{6}\right)
        \end{array}
        \longrightarrow
        \begin{pmatrix}1 & 0 & 0 & -\tfrac{4}{3}\\[4pt] 0 & 1 & 0 & \tfrac{4}{3}\\[4pt] 0 & 0 & 1 & \tfrac{1}{3}\end{pmatrix}
    \]

    Multipliziere Kennziffern und entnehmen:

    \[
        \begin{pmatrix}1 & 0 & 0 & -\tfrac{4}{3}\\[4pt]
                        0 & 1 & 0 & \tfrac{4}{3}\\[4pt]
                        0 & 0 & 1 & \tfrac{1}{3}\\[4pt]
                        0 & 0 & 0 & -1\end{pmatrix}
    \]

    \[
        \Rightarrow \text{für alle }
        \begin{pmatrix}b_1\\b_2\\a_1\\a_2\end{pmatrix}
        \in \left\{ \lambda \begin{pmatrix}-\tfrac{4}{3}\\[4pt]\tfrac{4}{3}\\[4pt]\tfrac{1}{3}\\[4pt]-1\end{pmatrix}
        \;\middle|\; \forall\lambda \in \mathbb{R} \right\}
    \]

    gilt die homogene Gleichung.

    Also:
    \[
        U \cap V
        = \left\{ \frac{\lambda}{3}\begin{pmatrix}1\\2\\3\end{pmatrix}
                  - \lambda\begin{pmatrix}-1\\2\\1\end{pmatrix}
          \;\middle|\; \forall\lambda \in \mathbb{R} \right\}
        = \left\{ \lambda\begin{pmatrix}1\\2\\3\end{pmatrix}
                  - 3\lambda\begin{pmatrix}-1\\2\\1\end{pmatrix}
          \;\middle|\; \forall\lambda \in \mathbb{R} \right\}
    \]
    \[
        = \left\{ \lambda \begin{pmatrix}4\\-4\\0\end{pmatrix}
          \;\middle|\; \forall\lambda \in \mathbb{R} \right\}
        = \left\{ -\frac{4}{3}\lambda\begin{pmatrix}-1\\0\\1\end{pmatrix}
                  + \frac{4}{3}\lambda\begin{pmatrix}0\\-1\\1\end{pmatrix}
          \;\middle|\; \forall\lambda \in \mathbb{R} \right\}
    \]

    Wegen $U \cap V = \left\{ \lambda \begin{pmatrix}4\\-4\\0\end{pmatrix} \;\middle|\; \forall\lambda \in \mathbb{R} \right\}$
    ist $\begin{pmatrix}4\\-4\\0\end{pmatrix}$ bzw.\ $\begin{pmatrix}1\\-1\\0\end{pmatrix}$ die Basis von $U \cap V$.

    \item \textbf{Bestimmung Basis von $U + V$:}

    Wegen a) und
    \[
        \dim(U + V) = \dim(U) + \dim(V) - \dim(U \cap V)
    \]
    folgt $\dim(U + V) = 3$. \quad $(\dim(U) = \dim(V) = 2)$

    Da $U + V$ ein Unterraum von $\mathbb{R}^3$ ist und
    $\dim(U + V) = \dim(\mathbb{R}^3)$ gilt, folgt $U + V = \mathbb{R}^3$.

    Die Einheitsvektoren $e_1, e_2, e_3$ sind also eine Basis von $U + V$.

\end{enumerate}

\end{document}
